\problemname{DetectiveDilemma}
% You are the new caretaker of an extremely weird child. Like all children, this child likes to ask lots of questions, questions as follow up to other questions, and even more follow up questions. This child however, has a very specific quirk: it can comprehend multiple explanations at the exact same time, meaning you can explain multiple related concepts to it at the exact same instance. However, if two things must be explained at the same time to make sense, and refer to each other, then explaining one thing before another will not make sense, and the child refuses to understand what you explain unless you explain them at the exact same time to it. Find the correct order to explain anything to the child. 

\Section*{Premise}
% \par{You are a bad detective trying to solve every case with as minimal (financial, physical and mental) average (of the three) effort as possible. The other detectives have worked out a connection between all the clues and suspects, but obtaining the clues/evidence is your job. You don't really care whether you convict the right person or not, so you just try to make the connection using minimal total effort to get every clue. To make it worse (for the city, better for you), the corrupt officials have put certain bribes in place, so some clues have "negative" financial burden for you (if you reach the same clue, they will not give you the bribe again). Rank the list of suspects by your standards of most to least suspicious, breaking ties by increasing alphabetical name order. }

\par{You are a bad detective trying to solve every case with \textit{minimal financial, mental and physical effort} as possible. The other detectives have worked out a connection between all the clues and suspects via a \textit{conspiracy-board} and it is your job to obtain the evidence and arrive at the \textit{suspect} to convict. You don't really care whether you convict the right person or not, as you try to make the connection from the \textit{initial clue} to any suspect spending the \textit{minimal total financial, mental and physical} effort as possible. Some clues connect only to other clues, while some clues connect to suspects. Connections between clues and suspects are considered "one way", as you don't really care about what happens if you have a suspect you can pin the crime on. To make it worse for the city, some corrupt officials have offered you a small "bonus" (a bribe) if you are to make certain connections between certain clues or suspects (instead of paying a financial cost, you get paid, or you pay a \textit{negative cost}). Of course, this isn't a way to make infinite money, as they won't offer you more money for arriving at the same conclusion twice (that connection won't incur any cost, but there is also no more bribe to be taken). Given your standards above, rank the list of suspects from most to least "suspicious". }

\Section*{Input}
\par{Line one consists of 3 numbers: $c\in\mathbb{N}$, the number of clues, $1\le c\le50$, $s\in\mathbb{N}$, the number of suspects $1\le s\e50$ and $n\in\mathbb{N}$, $0\le n\le c$, the nth clue that you start with. There are then $c$ more lines, given in the following format:

For each line, denoted line $C\le c$, the first number $0\le k\le c+s-1$ denotes the number the of outgoing connections to clues/suspects. There are then $k*4$ numbers, which can be grouped into groups of 4. The groups are as follows: the first number $0\le o\le c+s-1$ denotes the clue/suspect number that clue $C$ can point to. The next number $f\in\mathbb{Z}$ denotes the financial burden. $f\lt 0$ Implies a negative financial burden/bribe. The next two numbers $m, p\in\mathbb{N}$ are the mental and physical costs respectively. 

Note: The names of clues/suspects are zero-indexed, while the counts are all 1-indexed.
}

\Section*{Output}
\par{Output the most order of the suspiciousness of the suspects from most to least suspicious.}

% \par{You are the new caretaker of an extremely weird child. This child likes to ask a lot of questions and it is your job to answer them. However, the child is quite forgetful and can only remember the one topic if you explain every sub-topic related to that one topic during one explanation session. Furthermore, due to the forgetfulness of the child, the number of sub-topics are also the number of days for which the child will retain the information. You are quite lazy and decide to only explain topics for which the child will remember for the time you are spending with the child (you do not care if the child will ask the same questions after you leave). Explaining 1 topic, no matter the number of sub-topics will take exactly 1 day. Find the ordering of subtopics for which you will explain to the child.}

%\Section*{Input}
%\par{
%    The first line contains 2 naturals $c\in\mathbb{N}$, $\left(50\le c\le100,000 \right)$,$n\in\mathbb{N}$, $\left(c\le n\le 100\right)$
%}
% 1. Use Kosaraju's Algorithm to find the number of connected topics.
% 2. Rank the SCCs generated by 1 by size. to find the order. 
% 3. Keep only the first n SCCs and drop the rest, n being the number of days spent with the child. 


